Runtime monitors for long generations must warn before harm appears, yet activation readers such as probes and activation verbalizers answer only about the present or the next few tokens.
We ask whether an \emph{observer} with trajectory memory, learned future-window queries, a future-activation auxiliary loss and a censoring-aware hazard head forecasts further ahead than a matched direct readout of the latest activation.
We compare three such observers (horizon-query, latent rollout, hierarchical coarse-to-fine) against a parameter-, input- and epoch-matched direct readout on a frozen \qwen, with horizons of 64 to 2048 tokens and a prompt-family-level train/test split.
In a pilot of 84 on-policy trajectories, naturally occurring unsafe onsets are too rare (5 trajectories) and too early ($\le$448 tokens) to test long-horizon safety forecasting; a text-only TF-IDF baseline is best on this event.
On a denser ``caution'' event (first refusal or warning), the hierarchical observer reaches AUROC 0.976--0.990 versus 0.871--0.953 for the direct readout ($+0.04$ to $+0.11$), but family-bootstrap confidence intervals touch zero, the gain peaks at 256--512 tokens rather than growing with lead distance, and much of it reflects prompt-type separation.
Memory is the component with the clearest long-horizon benefit (removing it from the rollout observer drops AUROC at 1024 tokens from 0.944 to 0.831); the future-activation loss, hazard head and latent rollout show no benefit.
For long-form content, all observers forecast future-segment behavior at 0.83--0.88 accuracy up to 1024 tokens ahead and retrieve the remaining-response summary at 4--5$\times$ chance MRR.
Our results suggest that trajectory memory is worth its cost, and that testing long-horizon unsafe forecasting requires data with late natural onsets.
